\documentclass[11pt,a4paper,UTF8,fontset=fandol]{ctexart}

\usepackage{geometry}
\geometry{left=2.35cm,right=2.35cm,top=2.55cm,bottom=2.55cm}
\usepackage{graphicx}
\usepackage{booktabs,longtable,array}
\usepackage{amsmath,amssymb}
\usepackage{float,xcolor}
\usepackage[section]{placeins}
\usepackage{flafter}
\usepackage{hyperref,url}
\urlstyle{tt}
\makeatletter
\g@addto@macro\UrlBreaks{\do\a\do\b\do\c\do\d\do\e\do\f\do\g\do\h\do\i\do\j\do\k\do\l\do\m\do\n\do\o\do\p\do\q\do\r\do\s\do\t\do\u\do\v\do\w\do\x\do\y\do\z\do\0\do\1\do\2\do\3\do\4\do\5\do\6\do\7\do\8\do\9}
\setlength{\@fptop}{0pt}
\setlength{\@fpsep}{12pt}
\setlength{\@fpbot}{0pt plus 1fil}
\makeatother
\graphicspath{{assets/stage2_report_20261023/}}
\hypersetup{unicode=true,colorlinks=true,linkcolor=black,citecolor=black,
  urlcolor=blue,pdftitle={Stage2 RAW Video Denoising Research Report},
  pdfauthor={Stage2 DenoiseModel Project}}
\setlength{\parindent}{2em}
\setlength{\parskip}{0.30em}
\setlength{\emergencystretch}{3em}
\raggedbottom
\setcounter{topnumber}{3}
\setcounter{bottomnumber}{2}
\setcounter{totalnumber}{5}
\renewcommand{\topfraction}{0.92}
\renewcommand{\bottomfraction}{0.80}
\renewcommand{\textfraction}{0.08}
\renewcommand{\floatpagefraction}{0.85}
\setlength{\floatsep}{10pt plus 2pt minus 2pt}
\setlength{\textfloatsep}{12pt plus 2pt minus 2pt}
\setlength{\intextsep}{10pt plus 2pt minus 2pt}
\setlength{\abovecaptionskip}{6pt}
\setlength{\belowcaptionskip}{2pt}
\setlength{\LTpre}{6pt}
\setlength{\LTpost}{6pt}
\setlength{\tabcolsep}{4.5pt}
\renewcommand{\arraystretch}{1.16}
\newcolumntype{P}[1]{>{\raggedright\arraybackslash}p{#1}}
\newcommand{\dn}{\ensuremath{\,\mathrm{DN}}}
\newcommand{\db}{\ensuremath{\,\mathrm{dB}}}
\newcommand{\psnr}{\ensuremath{\mathrm{PSNR}}}
\newcommand{\rmse}{\ensuremath{\mathrm{RMSE}}}
\newcommand{\pairhead}{\makebox[0.49\textwidth]{输入}\hfill\makebox[0.49\textwidth]{方法3（EMA）}\\[0.25em]}
\newcommand{\triplehead}{\makebox[0.32\textwidth]{输入}\hfill\makebox[0.32\textwidth]{方法3}\hfill\makebox[0.32\textwidth]{方法8专用输出头}\\[0.25em]}
\newcommand{\djchead}{\noindent\makebox[0.25\textwidth]{\shortstack{含噪输入\\中心帧}}\makebox[0.25\textwidth]{\shortstack{官方 RViDeformer\\8帧}}\makebox[0.25\textwidth]{\shortstack{单帧基准\\1帧}}\makebox[0.25\textwidth]{\shortstack{五帧补偿\\第2300步}}\\[0.35em]}

\title{\textbf{视频去噪研究报告}\\[0.4em]
\large 方法探索、最终方案与多数据集泛化效果}
\author{肖纬杰\quad 邸锦承}
\date{}

\begin{document}
\maketitle
\begin{abstract}
\noindent
本报告记录 Stage2 RAW 视频去噪项目从公司数据适配、监督构建到跨传感器验证的研究过程。针对预训练教师在公司 RAW 上噪声与细节难以兼顾的问题，项目探索了两类方案：一类以 ExposureDiffusion 为基础，依次检查噪声建模、单帧适配、多帧与频域补偿、主干解冻，最后形成五帧输出残差候选；另一类直接微调 RViDeformer，并继续考察动态监督、传感器适配和推理融合。前一类方案虽降低了部分代理误差，但未取得稳定的细节恢复优势，未被采纳。最终采用方法3：以 RViDeformer-L 为基座，混合真实跨增益样本与标定噪声合成样本，使用带有效区域掩码的重建、梯度和高通损失进行全参数微调，并固定第5000步的 EMA 权重用于推理。sc635hai 另保留方法8的专用可选输出头。

效果展示覆盖公司 MIS20S1、SC465SL、sc635hai 和 sc450ai 多个数据集与拍摄条件。方法3在 MIS20S1 645x 跨增益参考评测中达到29.57\db，在未参与训练的1000x原生参考评测中达到53.27\db；在未进行目标传感器训练的 SC465SL 和 sc635hai 静态序列上，零样本 PSNR 分别提升约10.52--13.21\db 和13.59\db。报告同时展示高增益动态及 sc450ai 的实际输出，按静态定量、动态观察和专用适配三种协议解释结果。正文依次介绍方法探索、最终方案和多数据集效果，附录集中列出采纳方案的训练超参数、损失、模型与推理配置。

\vspace{0.6em}
\noindent\textbf{关键词：} RAW 视频去噪；RViDeformer；时域均值监督；跨传感器泛化；固定 ISP
\end{abstract}

\newpage
{\footnotesize\setlength{\parskip}{0pt}\setlength{\baselineskip}{11pt}\tableofcontents}
\newpage

\section{引言}

\subsection{项目起点与研究目标}

Stage1 使用预训练的 RViDeformer-L\_CRVD 教师生成伪干净目标。在公司 MIS20S1 数据上，该教师对图卡细线、棋盘和文字边缘的处理偏软，影响了下游监督质量。Stage2 因而研究如何得到更适合公司数据的 RAW 去噪模型：既尝试以另一种预训练恢复主干替代教师，也直接适配已有的视频教师，目标都是兼顾噪声抑制与高频结构，并检查不同传感器、增益和运动条件下的表现。

项目没有公司数据的逐像素干净动态真值。可利用的资源是低增益静态序列、同场景不同增益拍摄、平场与黑帧标定，以及少量其他传感器的真实场景流。因此，研究首先需要解决监督从哪里来，再判断相应监督能够支撑哪些模型改动。

\subsection{研究过程与报告主线}

报告先讨论 ExposureDiffusion 路线：公司域适配带来可测改善后，为什么仍需要引入邻帧；多帧模块确实训练后，为什么还要检查信息接入位置、主干冻结和监督方式；最后又如何根据固定验证与实际画面决定不采用这一候选。由此引出直接优化 RViDeformer 时空主干与监督的方案选择。

RViDeformer 路线中，方法0至方法4依次完成数据适配、参考构建、噪声合成、模型微调和评估，形成方法3。随后，动态内容被过度平滑的现象推动了方法5至方法7的运动保真探索；跨传感器测试引出了方法8至方法10的适配与退化建模；高增益动态预览又促成了方法11的置信门控实验。最终默认权重仍为方法3的第5000步 EMA，方法8仅保留 sc635hai 专用输出头。

两条路线的实际工作时间存在交叠，本文按问题与方案选择的逻辑组织，而非将章节顺序等同于实验日期。ExposureDiffusion 存档包含9月的噪声研究及10月18--25日的适配与补偿实验，方法3则已于9月23日完成验收。下文先写尝试与取舍，再详述采纳方法和多数据集效果；方法0至方法11沿用 RViDeformer 路线原编号，另一条路线用模型与实验名称区分。

\section{相关工作与数据基础}

\subsection{技术背景}

RAW 保留了传感器的线性响应和 Bayer 采样信息，适合在噪声统计层面建立恢复模型。CRVD/RViDeNet\cite{r1}、RViDeformer\cite{r2} 和 RVRT\cite{r3} 利用多帧信息进行恢复；HDR+\cite{r10}、Deep Burst Denoising\cite{r9} 与核预测方法\cite{r17} 则体现了对齐和帧间融合对低照度成像的作用。EMVD\cite{r15}、FastDVDNet\cite{r16} 以及 DnCNN\cite{r14}、CBDNet\cite{r20} 提供了不同的时空恢复与真实噪声建模思路；NL-means\cite{r18} 和 BM3D\cite{r19} 是非局部恢复的经典参照。本项目采用官方 CRVD 权重初始化 RViDeformer-L，在公司数据上微调。

在缺少干净真值时，Noise2Noise\cite{r4} 依赖噪声观测之间的统计关系，Noise2Void\cite{r5} 与 Noise2Self\cite{r6} 通过盲点构造自监督。静态长序列还可利用时域均值降低随机噪声，但运动区域的均值会混合不同时间的内容。项目据此将静态均值监督与动态探索分开处理，不把动态均值直接当作干净目标。

噪声标定以信号相关方差模型\cite{r7}为基础，并参考 RAW 数据合成的思路\cite{r8}。实际生成器采用逐增益经验系数、通道相关高斯采样、固定图样偏置和量化。动态探针使用 DIS 光流\cite{r11}。视觉比较在各自协议内冻结 ISP 参数；全数据集预览使用 Reinhard 色调映射\cite{r12}，奇偶参考面板和专用头比较沿用原审计渲染。

\subsection{公司数据集与实验分工}

全数据集预览覆盖4种传感器的45条原始 RAW 流，其中27条为场景流、18条为标定流。按传感器统计，MIS20S1 为29条（含18条标定），SC465SL 为12条，sc635hai 和 sc450ai 各2条。包含其他算法2DNR/3DNR处理的03组派生流不计入这45条。

表\ref{tab:dataset}采用 RViDeformer 主线的目录归类。ExposureDiffusion 交付包对其中10条型号未确认的465流保留 \texttt{465\_unspecified} 标识，按5个开发域统计；其训练、验证及测试划分也单独制定，见第\ref{sec:djc}节。同一条 RAW 在一条路线中留出，不意味着在另一条路线中也未参与训练。

\begin{table}[!htbp]
\centering\small
\caption{主要数据集、内容与实验用途。}
\label{tab:dataset}
\begin{tabular}{P{0.20\textwidth}P{0.29\textwidth}P{0.40\textwidth}}
\toprule
数据集 & 内容 & 用途 \\
\midrule
MIS20S1 暗房 & 同场景7档增益，10x--1000x & 10x静态主参考；128x/313x/632x/645x真实训练；1000x真实增益留出 \\
MIS20S1 室外 & 黄昏公园、人物运动，4档增益 & 场景迁移与高增益动态展示 \\
MIS20S1 标定 & 9档平场、9档黑帧，每流50帧 & 噪声拟合与合成验证；不作为场景泛化成绩 \\
SC465SL 暗室 & 256x、2000x，各300帧 & 未参与方法3训练的跨传感器静态定量 \\
SC465SL 超市与复杂纹理 & 增益梯子及 difficultpattern，共10流 & 多场景视觉与复杂纹理动态测试 \\
sc635hai & 131x低色温150帧、16x运动彩色200帧 & 方法3零样本测试；131x另用于方法8适配，16x保留 \\
sc450ai & 804x灯光变化150帧、1435x黑物运动210帧 & 跨传感器压力测试；1435x不参与适配 \\
\bottomrule
\end{tabular}
\end{table}

\subsection{RAW 适配与参考资格}

不同数据源的位对齐、CFA和黑电平并不统一。04场景流和05组通常需要从高12位右移4位；02组采用低位对齐。sc450ai 虽然附带说明写为高12位和4128\dn 黑电平，实测却是低位对齐、有效10位，样本均落在4\dn 格点，黑电平约256--258\dn。若直接套用说明，会把有效信号裁掉。场景流以 RGGB 为主，标定流为 GRBG；所有数据经适配器转换到统一的 R/Gr/Gb/B 通道顺序。

黑电平逐流、逐通道测量；帧数从文件字节数核实。标称增益不作为网络输入条件，因为不同档位伴随照明变化，实际噪声不随标签单调增加。例如暗房128x的分半绿色通道噪声约79.05\dn，1000x约28.93\dn。噪声合成因此使用标定实测系数，而非只按场景文件名推断强度。

监督资格随数据而定。10x低增益静态参考用于主方法训练；同流奇偶分半参考用于可信静态区域的测试或适配。旧教师输出、运动时域均值和盲点观测只作为对应探索的信号，不被登记为干净动态真值。

\subsection{评测口径}

静态评测在有效区域计算对应参考 DN 域的 PSNR、RMSE、高频保持率和残差。645x使用映射到10x亮度的参考域，1000x及跨传感器静态测试使用各自 native DN 域。跨传感器测试采用互斥奇偶帧池：偶数帧片段对奇数帧参考，再交换方向；只保留在全部、奇数、偶数帧池中均达到90\%静态置信度的像素。运动筛选采用 $\mathrm{median}+3\times1.4826\times\mathrm{MAD}$ 的鲁棒阈值。

动态序列没有干净逐帧目标，报告其时序响应、高频与裁剪等诊断量，并结合固定 ISP 视频观察。高频比下降既可能来自去噪，也可能来自结构损失；temporal MAD 下降既可能表示抑制噪声，也可能表示运动被平均。因此，两者均与视觉结果一起解释。

所有同序列输入与输出共享白平衡和曝光，关闭 ISP 的2DNR/3DNR。全数据集展示采用 edge-aware 去马赛克、高光重建和 Reinhard 映射。PSNR 在 RAW 域计算，展示视频用于观察降噪、纹理与运动；不同数据集的参考域和有效区域不同，绝对 PSNR 不作跨数据集排名。

以上静态参考指标用于 RViDeformer 主线。ExposureDiffusion 路线主要报告静态有噪参考与运动教师目标构成的代理误差，以及同帧固定显示参数下的视觉结果。该误差没有干净真值含义，不能换算为本报告静态参考表中的 PSNR，也不能把不同裁块、掩码和损失定义下的百分比相加。

\section{ExposureDiffusion与多帧补偿的探索}
\label{sec:djc}

\subsection{起点：先让噪声模型与公司RAW对应起来}

这条路线尝试直接从含噪 RAW 得到去噪 RAW，以公司适配模型承担恢复任务，推理不依赖公司2DNR或3DNR输出。起初的问题是：除模型预训练域不同外，公司 RAW 的噪声、量化与固定图样是否也没有被正确表达？因此，存档中的早期工作先从噪声生成和观测模型展开。

单暗帧频谱方案以每档第一张暗帧建模，其余49张作检查。平均频谱与总体噪声量级接近真实数据，但高增益的稳定高频固定图样被部分随机化，合成的时间噪声标准差高出约4\%--7\%。进一步检查还发现，暗帧存在2\dn{}码格及逐帧奇偶状态变化；此前按1\dn{}取整的分布距离不能继续使用。研究由此从“匹配一张暗帧的分布”推进到“区分固定图样、帧级偏置与随机残差”。

随后使用40张暗帧拟合多帧分解，并与剩余10张作同条件比较。以25600档为例，固定图样 RMSE 从单帧方案的9.628\dn{}降到6.858\dn，时间标准差相对误差从$+6.63\%$降到$+0.08\%$；但合并码值分布的 Wasserstein 距离从0.194\dn{}增到0.236\dn。低增益也未表现出全面优势。这个结果促使评价从单一平均频谱转向固定结构、时间方差和码值分布的共同检查，不能仅凭一个指标决定噪声模型已经适用。

散粒噪声拟合也经历了修正：最初把约83毫秒暗帧的方差固定为不同曝光彩色帧的截距，后来改用彩色重复帧对拟合，并在内部选择暗截距、非负截距或零截距。25600档的末5对检查误差由19.74\%降至9.85\%。不过，参与该检查的最低扣黑信号仍为48\dn，高增益系数保留为暂定估计，不能直接把结果外推到近黑区。

存档还记录了泊松散粒项、高斯混合暗残差与量化区间概率组成的可微似然，以及静态背景的条件残差模型。量化似然修正了正态分布函数的重复缩放，并重新进行了九档拟合与暗帧留出检查；条件残差候选在同记录的时间、位置留出中得到最差逐通道方差相对误差3.12\%。这些工作提供了噪声统计证据，但暗场曝光与有光场景并不完全一致，尚不能据此宣称真实场景后验重建成功。七帧联合似然、DPS与条件更新代码亦被保留，交付记录未提供可对应的完整七帧扩散去噪训练及效果结果。实际模型适配于是采用可核查的真实帧监督，未将这些候选合成器直接登记为最终训练退化。早期噪声记录与后续训练产物的来源分别列于附录\ref{app:djc_sources}。

进入多传感器开发后，又逐流检查位对齐、CFA、黑电平与信号相关方差。一次具体修正发生在运动掩码：若先把扣黑均值截为非负，黑电平增加8\dn{}会使最暗465流约26.96\%的位置改判为静态；保留有符号统计均值、仅给最终方差设置正下限后，复核的最大分类变化降至0.4349\%。只有先把这类统计偏差厘清，静态有噪帧对才具有训练意义。

\subsection{单帧适配：误差下降后，细节问题仍然存在}

模型采用 ExposureDiffusion 的官方 SID 预训练 NAFNet 权重。主干包含四级编码器、12个中间块和四级解码器，通道宽度从16增加到256再恢复到16。每次前向接收4通道 RAW 状态及4通道光子条件，同一网络共享参数进行三次曝光递推，并保留泊松采样。这里的“三次”是递推次数，不是三帧输入；虚拟曝光比例100也不是传感器增益标签。

本路线的45条流被分为20条开发流、18条标定流和7条 MIS20S1 暗房测试流。开发流包含2条 SC450AI、2条 SC635HAI、2条 SC465SL、4条 MIS20S1 室外流及10条型号未确认的465流。每条开发流的前约70\%用于训练，间隔10帧后使用尾段验证；输入、参考、引导与教师窗口均须完整落在所属区间。7条暗房流不参与本路线的训练和权重选择，属于暗房序列留出；由于开发集也有 MIS20S1，且其他传感器可能拍摄相同室内布景，不能称为传感器型号与物理布景均未见过。

中心输入为$x_t$，静态区使用独立有噪帧$x_{t+6}$作为目标；运动区使用冻结官方 CRVD 权重的八帧 RViDeformer 输出$T_t$作为软目标。定义$M_s$和$M_m$为静态、运动掩码，单次输出$z$的训练误差为
\begin{equation}
 \mathcal L_{\mathrm{mix}}(z)=
 \operatorname{MSE}_{M_s}(z,x_{t+6})
 +0.5\operatorname{MSE}_{M_m}(z,T_t).
 \label{eq:djc_mix}
\end{equation}
式中$x$、$T$与$z$均为扣除黑电平后按逐通道有效范围$4095-b_c$归一化的量，因此代理 MSE 为无量纲数值。两类误差分别按各自有效元素数归一化，空区域不贡献相应项；三次递推输出的损失再平均。七张独立引导帧只用于构造掩码，并非七帧网络输入。噪声归一化的时域变化用于区分静态、运动与忽略区，过亮及不可靠区域剔除。教师只参与训练监督和对照，推理无需教师。

单帧主干全参数适配1000步，使用FP32、$128\times128$ packed裁块、batch size 32、AdamW、学习率$10^{-5}$、零权重衰减及梯度裁剪1。原固定40裁块的代理误差从$1.802499\times10^{-4}$降至$1.602098\times10^{-4}$，下降11.118\%，最优权重在第1000步。公司域适配确实有效，但128x留出视频多个时刻的文字、线条和图卡仍偏软。这使下一轮实验不再只增加单帧训练，而是尝试让邻帧提供中心帧中较弱的结构信息。

\subsection{从五帧残差到特征与频域补偿}

第一轮多帧尝试冻结公司单帧主干，增加 BIPNet 五帧 RAW 残差分支，用$t-2$到$t+2$补偿最终输出。新增分支约667万参数，从头训练250步。固定代理误差仅从$3.254105\times10^{-4}$降至$3.253263\times10^{-4}$，下降0.0259\%。进一步将五帧替换为重复中心帧，第60帧输出的差异仅0.002104\dn；多组画面中，残差能量的74.8\%--86.0\%可由常量通道偏移解释。网络虽产生了非零修正，却没有给出文字和分辨率卡恢复的明确证据。

这一诊断把问题从“是否需要更多参数”推进到“邻帧信息究竟在哪里进入、学到的是否为内容”。因此，后续分别尝试在第一次曝光递推中注入 TAP 对齐特征、在特征域加入 AdaIR 频域补偿，以及在输出域加入 ORNet 残差。TAP 从冻结主干抽取五帧64通道特征，经金字塔可变形对齐与融合后注入对应解码跳接；AdaIR和ORNet这两项实际使用单帧信号，作为不同补偿机制的对照，不能都算作时序模型。

\begin{table}[!htbp]
\centering\small
\caption{不同补偿位置的实际试验。前三项共用混合代理误差基线$3.254105\times10^{-4}$；多尺度首步另用独立评测设置。}
\label{tab:djc_pilots}
\begin{tabular}{P{0.21\textwidth}P{0.12\textwidth}P{0.22\textwidth}P{0.33\textwidth}}
\toprule
尝试 & 实际步数 & 最优代理误差 & 实验后的判断 \\
\midrule
TAP首步特征注入 & 750 & $3.190954\times10^{-4}$，第250步 & 下降1.941\%，邻帧作用仍弱 \\
AdaIR频域补偿 & 750 & $3.157863\times10^{-4}$，第750步 & 下降2.958\%，视觉细节未形成稳定优势 \\
ORNet输出残差 & 750 & $3.214685\times10^{-4}$，第750步 & 下降1.211\%，与原输出总体接近 \\
确定性首步、三级TAP & 750 & 单看MSE第250步略降 & 梯度项未通过，正式选择退回第0步 \\
\bottomrule
\end{tabular}
\end{table}

前三项在补充裁块中仍有小幅代理误差下降，但图卡和字牌的改善并不与数值同比例出现。逐递推诊断还发现，一个固定设置下第三步代理误差高于第一步，且泊松路径会影响比较。因此又尝试确定性首步恢复，并按64、32、16通道逐级增加 TAP 融合，各训练250步。第250步的 MSE 仅略降，梯度误差却升高；第750步 MSE 也回退。按“重建与梯度均不劣于基线”的规则，正式选择仍是未训练的第0步。

同时进行的无训练对照，在静态掩码内平均五帧 RAW 或五次去噪输出。在一个首步设置下，先平均五帧RAW、再经冻结的单步 ExposureDiffusion 恢复，使代理 MSE 从$2.62842\times10^{-4}$降至$1.68969\times10^{-4}$；但完整三步的另一设置没有同样结论。它说明重复观测中存在可利用信息，也说明简单平均及按帧数缩放散粒噪声不足以替代对运动、读出噪声与固定图样的处理。

\subsection{进一步检查监督与主干冻结的影响}

既然附加模块的收益较小，研究继续检查两种可能：冻结主干是否限制了新特征的使用；运动教师目标是否使模型更倾向于复现已有平滑结果。相应实验取消教师监督，只在可信静态区域使用真实帧对，并依次解冻解码器和完整主干。训练目标改为$t+4,t+6,t+8$三帧平均，验证仍使用$t+6$；评价采用静态 MSE、细节 MSE 与0.25倍梯度 MSE 之和。

\begin{table}[!htbp]
\centering\small
\caption{静态监督与主干更新实验。各行按自身起点比较，指标是组合目标，不是干净真值PSNR。}
\label{tab:djc_unfreeze}
\begin{tabular}{P{0.22\textwidth}P{0.23\textwidth}P{0.28\textwidth}P{0.14\textwidth}}
\toprule
尝试 & 训练与继承 & 起点$\rightarrow$最优目标 & 相对下降 \\
\midrule
解码器与三级TAP联合训练 & 2500步，最佳2250步 & $7.850996\times10^{-4}$ $\rightarrow 7.588867\times10^{-4}$ & 3.339\% \\
完整主干继续微调 & 继承上行最佳，再2500步 & $7.588867\times10^{-4}$ $\rightarrow 7.446752\times10^{-4}$ & 1.873\% \\
官方TAP64公司适配 & 独立预训练初始化，1500步 & $7.635421\times10^{-4}$ $\rightarrow 7.352122\times10^{-4}$ & 3.710\% \\
\bottomrule
\end{tabular}
\end{table}

联合解码器训练后，真实五帧与重复中心帧的验证输出差异达到约9.31\dn，说明模型开始明显使用邻帧。然而，局部玩偶与纹理仍偏软，输出差异本身不能证明真实结构已经恢复。独立的官方 TAP64 预训练对照在部分局部更清楚，也提示预训练表征和完整视频主干值得重视。上述结果没有把 ExposureDiffusion 改造方案推到可稳定替代原教师的程度，因此保留为探索记录。

\subsection{回到清晰的输出残差：五帧候选的完整验证}

后续实现曾组合 TAP 对齐、DarkIR模块与 FrENet 频域模块，通过内部钩子注入特征。服务器准备快照记录的仍是未更新状态，迁移说明另记录该组合版本曾试训20步；这些机械检查不能计作有效效果结论。最终候选在交付代码中标为V3，删除内部注入与复杂组合，恢复完整冻结的三步 ExposureDiffusion，在最终 RAW 输出处叠加一条五帧残差分支：
\begin{equation}
 \widehat y_t=B(x_t)+R_\theta(x_{t-2:t+2},\eta),
 \label{eq:djc_v3}
\end{equation}
其中$B$为公司适配单帧主干，$\eta$为噪声条件。残差分支共享卷积提取16、32、64通道的三级特征，用 TAP 的金字塔对齐与融合组件汇聚五帧，再生成4通道 RAW 修正。输出层零初始化，使训练起点等于冻结主干；分支约70万参数，全模型约760万参数。它不把固定高通结果强行加回画面，推理也不需要教师或参考帧。

这一版本继续使用式\eqref{eq:djc_mix}，只约束最终输出。五帧输入与$t+6$目标互不重叠，8张独立引导帧负责掩码，静态区域还需通过更保守的局部一致性检查。实际训练为FP32、$192\times192$ packed裁块、batch size 2、梯度累积4、有效batch 8；AdamW学习率$10^{-4}$、权重衰减0、梯度裁剪1，共5000次优化器更新。固定验证每条开发流取1个裁块，共20个，先在域内汇总，再对5个开发域等权平均，每100步验证一次。

\begin{figure}[!htbp]
\centering
\includegraphics[width=\textwidth]{djc/djc_training_curves.pdf}
\caption{单帧适配与五帧输出补偿的固定验证曲线。左图使用单帧阶段的40裁块；右图使用五帧阶段的20裁块及等域汇总。两图各自比较训练前后的代理误差，不能直接比较绝对值或累加下降比例。五帧候选在第2300步最佳，继续训练至5000步后反而高于冻结基线。}
\label{fig:djc_training}
\end{figure}

如图\ref{fig:djc_training}，五帧阶段的冻结基线为$1.794418\times10^{-4}$，第2300步降至$1.772457\times10^{-4}$，下降1.224\%；第5000步升至$1.804805\times10^{-4}$，比基线高0.579\%。因此，在额外测试前已固定第2300步为本路线候选，第5000步仅作训练终点对照。它与单帧适配的11.118\%使用不同验证设置，不能合成为总改善比例。

\begin{figure}[!htbp]
\centering
\includegraphics[width=0.94\textwidth]{djc/djc_domain_results.pdf}
\caption{五帧候选的分域固定验证。正值表示代理误差相对冻结单帧基准下降。所有域均在开发训练中出现，属于同域时间留出；465型号未确认的10条流独立统计。改善随数据域和训练阶段改变，没有形成各域一致的优势。}
\label{fig:djc_domains}
\end{figure}

图\ref{fig:djc_domains}的分域结果进一步解释了为什么总体数字不足以支持采用。第2300步在 SC465SL 和 SC635HAI 分别下降5.477\%和2.229\%，MIS20S1却增加0.613\%；到第5000步，SC465SL又变为增加6.694\%。复查固定裁块发现，SC465SL其中一个裁块只有2个有效运动位置，SC635HAI其中一个只有1个。由于静态和运动误差分别归一化，这些少量位置仍承担0.5倍运动项权重，代理误差会对局部教师目标很敏感。这里的位置指 packed Bayer 空间位置，并在四通道上计算损失。

补充实验在相同开发视频中另取两组20裁块，第2300步分别改善0.352\%和0.931\%，第5000步反而改善0.807\%和2.283\%。这些是裁块位置种子的变化，不是新数据集或独立训练重复，也没有据此重选权重。邻帧诊断中，第2300步将真实五帧换为重复中心帧，输出差异RMS的中位数仅0.0134\dn；反转邻帧顺序为0.00834\dn，而残差RMS的中位数为0.880\dn。三项均先在显示ISP前的12位DN域逐裁块计算RMS，再取20裁块的中位数。综合起来，候选确有小幅统计收益，但其排序受裁块构成影响，最佳候选的邻帧作用也较弱，仍需回到实际画面判断结构改善。

\subsection{从实际画面决定是否采用}

交付结果覆盖27条场景流的165个同帧案例：20条开发流各3帧，共60例；7条暗房留出流各15帧，共105例。每例保留输入、官方 RViDeformer、公司单帧基准、第2300步候选和第5000步终点。全图与局部面板只是同一案例的不同展示，不另计样本数。下列四列均使用逐流固定的白平衡与亮度映射，展示参数从该流首个展示输入帧确定；它们没有干净真值。

\begin{figure}[!htbp]
\centering\small
\djchead
\includegraphics[width=\textwidth,trim=0 0 0 112bp,clip]{djc/djc_1000x_full.png}\\[0.35em]
\includegraphics[width=\textwidth,trim=0 0 0 112bp,clip]{djc/djc_1000x_detail.png}
\caption{MIS20S1暗房1000x序列第99帧：全图与同位置局部。该序列未参加本路线训练及选权重。四列依次为输入、官方CRVD权重的八帧 RViDeformer、公司单帧基准、五帧输出补偿第2300步。后两列总体接近，棋盘格和玩偶衣物条纹在官方 RViDeformer 输出中更易辨认。此处的 RViDeformer 是冻结官方参照，并非后文方法3微调权重。}
\label{fig:djc_1000}
\end{figure}

图\ref{fig:djc_1000}体现了这条路线最直接的取舍依据：单帧适配与残差补偿都能压低噪声，但在高噪声图卡和玩偶区域，五帧候选仍接近单帧基准，未显示出足够明确的细节恢复。官方 RViDeformer 虽然也存在平滑问题，在此样本中仍保留了更可辨的结构。这不是等参数或等帧数消融，而是项目能否用当前候选替换现有视频教师的实际对照。

\begin{figure}[!htbp]
\centering\small
\djchead
\includegraphics[width=\textwidth,trim=0 0 0 112bp,clip]{djc/djc_10x_detail.png}
\caption{MIS20S1暗房10x序列第8帧，列顺序同图\ref{fig:djc_1000}。该流亦为本路线留出。低噪输入原本含有较清楚的结构，模型间差异小于1000x样本，五帧补偿仍未呈现普遍、清晰的细节增益。该流在后文方法3中用于构建参考，数据身份随路线分别记录。}
\label{fig:djc_10}
\end{figure}

低噪留出样本和开发域尾段用于检查上述判断是否只来自最难样本。图\ref{fig:djc_10}中差别整体较小；图\ref{fig:djc_sensors}中，SC450AI人物与SC635HAI图卡的五帧输出仍与单帧基准接近。分域代理误差下降并不意味着每个实际局部都获得可见收益。这批静帧也不提供连续视频闪烁或重影改善的结论。

\begin{figure}[!htbp]
\centering\small
\djchead
\includegraphics[width=\textwidth,trim=0 0 0 112bp,clip]{djc/djc_sc450_detail.png}\\[0.4em]
\includegraphics[width=\textwidth,trim=0 0 0 112bp,clip]{djc/djc_sc635_detail.png}
\caption{开发域的时间留出画面：上为SC450AI灯光变化804x序列第132帧，下为SC635HAI低色温131x序列第132帧。四列仍为输入、官方 RViDeformer、单帧基准和五帧第2300步候选。两流的前段参与训练，图示尾段属于验证，不能称为新传感器零样本结果。}
\label{fig:djc_sensors}
\end{figure}

\subsection{未采纳的原因与RViDeformer路线的衔接}

这条探索并非所有环节都无效：公司单帧适配明显降低了原固定代理误差，部分补偿和解冻实验也产生了可测收益。但项目需要的是可稳定替换现有教师的 RAW 视频恢复结果。经过多种信息接入位置、监督方式和更新范围的尝试，最后五帧候选的整体增益仍小，随域和裁块变化，真实图卡与纹理也没有建立稳定优势。因此，在本项目已有数据和已完成实验范围内，这条改造路线未达到采用要求，最终不作为 Stage2 默认方案。

这些实验给方案选择带来的启发是明确的：单纯增加外围补偿模块，并不能保证网络利用邻帧恢复真实结构；而在实际对照中，已有视频教师仍体现出可用的时空表征。报告由此转入另一条路线：直接把 RViDeformer 本身作为优化对象，同时用可信的静态时域参考和实测噪声构造监督，让公司域适配发生在完整时空主干内。下一节说明这一方案如何形成方法3，以及为什么后续尝试最终仍保留它。

\section{RViDeformer微调与后续方案选择}
\label{sec:exploration}

\subsection{方法0至方法4：形成静态基线}

方法0建立逐流适配、配准、曝光映射与运动掩码。方法1用暗房10x的200帧构建静态参考，奇偶均值残差为0.51--0.80\dn，曝光归一化后的跨增益偏差为1.17--1.86\dn。方法2利用平场和黑帧生成9档合成噪声，最大逐通道统计误差8.66\%，低于15\%验收阈值。上述准备共同支撑方法3的混合数据微调，方法4负责评估。

评估中出现过一次口径修正：1000x与10x跨捕获曝光关系在高亮处偏离线性，同一模型按该映射仅有18.84\db；配准复核后，改用1000x自身的静态参考得到53.27\db，奇偶互斥验证为53.10/53.07\db。该结果使项目保留方法3，并将跨捕获映射与模型误差分开检查。方法3的监督、损失和推理流程在第\ref{sec:adopted}节统一说明。

\subsection{方法5至方法7：运动保真探索}

方法3的静态去噪改善后，645x视频中的运动人物仍有半透明残影。原运动审计中，输出的运动区帧差仅为输入的27\%--32\%。这推动了三条后续路线，比较时均要求运动响应改善且静态、原生增益及高频指标不回退。

\begin{table}[!htbp]
\centering\small
\caption{运动保真路线的简要比较。}
\label{tab:motion_trials}
\begin{tabular}{P{0.10\textwidth}P{0.30\textwidth}P{0.38\textwidth}P{0.10\textwidth}}
\toprule
方法 & 设计 & 主要结果 & 取舍 \\
\midrule
方法5 & 光流对齐的动态 N2N；教师静态锚；冻结范围与损失权重实验 & 收紧运动置信规则后覆盖仅1.81\%--3.82\%；多组候选未同时通过运动、静态和1000x门槛 & 未采纳 \\
方法6 & 冻结方法3，训练空间残差运动专家；共4版 & 门控 IoU 达0.533，但真实动态内容未稳定恢复；更强残差伴随静态偏移或目标误差增大 & 未采纳 \\
方法7 & 在运动区融合旧教师残差，扫描36组配置 & 最强配置的运动比相对增加0.9162\%；固定 ISP 未见稳定的人物轮廓改善 & 未采纳 \\
\bottomrule
\end{tabular}
\end{table}

方法5用 DIS 光流将后帧对齐到当前帧，以前后向误差和光度误差筛选可信像素。增加严格的1像素运动门槛后，监督覆盖明显缩小。方法6转而冻结主干，减少全模型漂移；方法7再取消训练，检验旧教师残差是否能直接补充运动信息。三者依次减少了参数更新带来的影响，但在既定数据与验收条件下，均未得到足以替换基线的结果。

因此，后续实验继续保留方法3，研究重心转向跨传感器：先测同一权重在哪些数据域有效，再决定是否需要目标域适配。

\subsection{方法8至方法10：跨传感器适配与退化验证}

零样本测试表明，方法3在 SC465SL 暗室和 sc635hai 低色温静态域有明显降噪收益，而 sc450ai 黑色运动物体仍明显发软。方法8据此考察逐传感器适配：sc635hai 使用奇偶静态监督，sc450ai 使用单帧盲点监督，均混入 MIS20S1 回放。结果仅保留 sc635hai 的专用输出头，详细配置见第\ref{sec:head}节；sc450ai 的5个检查点均没有恢复高频，1435x端点裁剪还从3.28\%升至4.64\%。

方法9尝试统一微调4种传感器，只更新15,508个参数、训练100步。SC465SL和sc635hai静态指标略升，但 sc450ai 的高频比由0.126621降至0.126589，未改善黑衣与运动轮廓，故没有作为统一模型采纳。

方法10随后检查 sc450ai 的监督来源。现有数据只有两条动态场景流，没有独立暗场、平场或干净 RAW 配对。其补充退化实验比较了原独立噪声、4\dn{}量化独立噪声和频谱整形噪声；频谱整形虽改善804x的聚合相关误差，却未通过逐通道噪声与1435x保留集验证。该生成器未进入最终方法。

这组实验使最终方案保持清晰：默认使用同一方法3权重评测泛化，仅给 sc635hai 保留一个专用可选头。目标域的零样本结果与专用适配结果分别报告。

\subsection{方法11：高增益推理门控}

全数据集预览显示，MIS20S1室外1500x等序列仍有运动朦胧。四路诊断比较原始输入、正常时域推理、加权空间拼接和单帧重复8次的推理：改变拼接权重对主体影响很小，单帧重复则恢复了更多当前帧结构，同时带回明显彩噪。这些观察支持将时域融合视为该现象的重要来源。

方法11尝试用输入残差、梯度和光流置信度，在正常时域输出与单帧重复输出之间融合。两条代表序列的非零门控覆盖仅0.487\%和0.713\%，全局输出接近原时域分支，逐帧观察也未见稳定的主体重影改善，因此未纳入最终推理流程。

\subsection{最终方案的确定}

从 ExposureDiffusion 替代方案到 RViDeformer 直接微调，再到运动与传感器专项实验，研究逐步区分了代理误差下降、实际结构恢复与跨域可用性。最终保留的方法3具有完整的监督与评测记录，静态恢复结果稳定；其余默认替代候选未在各自目标上满足采用条件。方法8的 sc635hai 输出头以专用可选权重单独保留。接下来围绕这两项说明最终实现，并将其效果统一放到多数据集测试中展示。

\section{最终采纳的方法}
\label{sec:adopted}

\subsection{整体流程}

默认方法由四个环节组成：RAW适配、可信静态监督、真实与合成混合训练、EMA多帧推理。方法0至方法2提供的数据处理和监督构造均属于方法3的完整训练管线。

\begin{figure}[!htbp]
\centering\small
\fbox{\parbox[c][1.5cm][c]{0.25\textwidth}{\centering 原始 RAW\\位对齐、CFA、黑电平\\配准与曝光映射}}
$\longrightarrow$
\fbox{\parbox[c][1.5cm][c]{0.25\textwidth}{\centering 静态参考与有效掩码\\实测噪声合成\\真实:合成约1:1}}
$\longrightarrow$
\fbox{\parbox[c][1.5cm][c]{0.25\textwidth}{\centering RViDeformer-L\\全参数微调5000步\\EMA权重推理}}
\caption{最终方法3的完整流程。训练数据进入统一参考域；输出在对应协议的参考 DN 域评测，再使用该协议的固定 ISP 展示。}
\label{fig:adopted_pipeline}
\end{figure}

模型输入为8帧、4通道的 packed Bayer，输出对应的8帧。训练采用 $64\times64$ packed patch，即原 Bayer 的 $128\times128$ 区域；推理通过重叠分块处理完整画面。方法3不增加增益标签输入或传感器条件分支。

\subsection{网络架构与数据流}
\label{sec:network_flow}

方法3沿用 RViDeformer-L 主干\cite{r2}，整体见图\ref{fig:rvideformer_overview}。输入先由浅层卷积从4通道映射到84通道，再进入 U 形时空重建模块。七个阶段相对 packed RAW 的空间尺度依次为$1,\,1/2,\,1/4,\,1/8,\,1/4,\,1/2,\,1$，时间长度保持8帧；同尺度编码与解码特征逐元素相加，将粗尺度上下文与细尺度结构结合。

\begin{figure}[!htbp]
\centering
\includegraphics[width=\textwidth]{assets/RVIDeformer/屏幕截图 2026-01-30 151658.png}
\caption{RViDeformer整体架构\cite{r2}。左侧为U形时空重建，右侧为全分辨率特征细化，长跳接用于残差重建。入口Conv2d在本项目中对应时间核为1的$1\times3\times3$卷积，逐帧提取空间特征。}
\label{fig:rvideformer_overview}
\end{figure}

\textbf{空间信息如何聚合。}图\ref{fig:rvideformer_mssa}中的 MSSA 以当前窗口生成查询$Q$，由局部窗口、其低分辨率版本及更大邻域或全局缩略特征提供键$K$和值$V$。各分支按相似性加权聚合，再拼接并映射融合；移窗操作让相邻窗口交换信息，从而同时利用局部细节与较大范围的相关结构。

\textbf{邻帧信息如何进入。}图\ref{fig:rvideformer_mtsa}中的 MTSA 以参考帧提供$Q$，从支持帧的对应窗口及更大邻域或全局范围提供$K,V$，通过注意力建立隐式帧间对应。跨帧聚合结果再与 MSSA 分支融合；前馈层和局部卷积继续更新特征，残差连接保留原有信息。图中的 RepMLP、RepConv 为原论文的可重参数化设计，本文使用其训练形式。

\textbf{特征如何还原为RAW。}U形模块回到全分辨率后，将84维特征映射到108维，经四组细化模块再映射回84维，并与浅层特征相加。末端卷积生成4通道残差，与输入相加，得到对应的8帧去噪RAW。本文细化段的自注意力窗口仍可跨帧聚合；具体深度、窗口与7,505,716个可训练参数的配置见表\ref{tab:model_config}。本项目在该主干上改变公司数据监督与微调方案。

\begin{figure}[H]
\centering
\includegraphics[width=0.82\textwidth]{assets/RVIDeformer/屏幕截图 2026-01-30 183700.png}
\caption{多分支空间自注意力MSSA\cite{r2}。不同分支从不同空间范围提取信息；图中C表示拼接，圆圈加号表示逐元素相加。}
\label{fig:rvideformer_mssa}
\vspace{0.6em}
\includegraphics[width=0.82\textwidth]{assets/RVIDeformer/屏幕截图 2026-01-30 204604.png}
\caption{多分支时域互注意力MTSA\cite{r2}。参考帧生成查询，支持帧提供键和值，融合局部与较大范围的跨帧信息。}
\label{fig:rvideformer_mtsa}
\end{figure}

\FloatBarrier
\subsection{真实样本与静态监督构建}

主参考来自 MIS20S1 暗房10x序列的200帧平均。参考构建包含固定图样修正，并记录运动及无效区域。参考形状为 $4\times540\times960$。奇偶分半残差在 R/Gr/Gb/B 上分别为0.511、0.797、0.763和0.560\dn；图卡高频检查确认平均后仍保留可用结构。

真实输入选用128x、313x、632x和645x。对每个输入增益，先读取其位对齐、黑电平和运动记录，再使用整数配准及逐通道仿射曝光映射进入主参考坐标和亮度域。设映射后的像素为 $x_c$，参考黑电平为 $b_c$，归一化为
\begin{equation}
 \widetilde{x}_c=\operatorname{clip}\left(\frac{x_c-b_c}{4095-b_c},0,1\right).
 \label{eq:normalization}
\end{equation}
目标为相同位置的静态参考，在8个时间位置上重复。

有效掩码同时满足参考有效、输入8帧均为静态、全部输入通道均未饱和。饱和阈值为映射域3800\dn，运动掩码阈值为0.5。只有这些位置参与损失，避免对错位运动和饱和内容施加静态均值目标。1000x真实序列不进入训练。

\subsection{实测噪声合成与混合采样}

合成分支仍以同一静态参考为干净目标，通过9档平场与黑帧建立噪声模型。信号相关方差可写为
\begin{equation}
 \sigma_c^2(s,G)=a_c\,sG+b_c^{\mathrm{read}}G^2.
\end{equation}
实际生成器保留了每档增益的经验 shot 系数 $\alpha_{c,g}$ 和 read variance $v_{c,g}$，用于处理标定中的非单调变化。令 $d$ 为固定图样偏置，$z\sim\mathcal N(0,I)$，$LL^\top=C$ 为聚合通道相关矩阵的 Cholesky 分解，则合成为
\begin{align}
 q &= y+d,\qquad s_c=\max(q_c-b_c,0),\\
 x^{\mathrm{syn}} &= \operatorname{clip}\left(
 \operatorname{round}\left[q+\sqrt{\alpha_g\odot s+v_g}\odot(Lz)\right],0,4095\right).
 \label{eq:synthetic}
\end{align}
该实现使用相关高斯采样匹配信号相关方差。合成后按1\dn{}格点量化，保留原始黑电平与合法范围。

合成增益标签为100、200、400、800、1600、3200、6400、12800和25600，每档生成8个独立噪声时序变体。8帧具有相同干净内容和不同噪声，真实与合成分支按约1:1混合。真实性检验比较合成与真实同档数据的高通差分噪声，最大逐通道相对误差8.66\%，通过15\%阈值。逐档参数列于附录。

\subsection{优化目标}

方法3严格加载官方 CRVD 权重后进行全参数微调，网络结构与数据流见第\ref{sec:network_flow}节。

训练先将预测裁剪到 $[0,1]$，再计算带掩码的损失。记
\begin{equation}
 \langle z\rangle_M=\frac{\sum M\odot z}{\max(\sum M,1)},
\end{equation}
其中掩码沿通道维展开。总损失为
\begin{equation}
 \mathcal L=\mathcal L_{\mathrm{char}}+0.1\mathcal L_{\mathrm{grad}}+0.1\mathcal L_{\mathrm{hp}}.
 \label{eq:loss}
\end{equation}
Charbonnier 项\cite{r13}约束重建误差：
\begin{equation}
 \mathcal L_{\mathrm{char}}=\left\langle\sqrt{(\hat y-y)^2+10^{-6}}\right\rangle_M.
\end{equation}
若写为 $\epsilon^2$，则 $\epsilon=10^{-3}$。梯度项约束水平、垂直梯度的幅值差：
\begin{equation}
 \mathcal L_{\mathrm{grad}}=\tfrac12\left(
 \left\langle\left||D_x\hat y|-|D_xy|\right|\right\rangle_{M_x}+
 \left\langle\left||D_y\hat y|-|D_yy|\right|\right\rangle_{M_y}\right),
\end{equation}
其中 $M_x,M_y$ 是相邻有效掩码的乘积。高通项为
\begin{equation}
 H(u)=u-\operatorname{AvgPool}_{3\times3}(u),\qquad
 \mathcal L_{\mathrm{hp}}=\langle|H(\hat y)-H(y)|\rangle_M.
\end{equation}
高通仅在空间计算，不跨时间平均。重建项负责整体恢复，两个0.1权重的结构项补充边缘和局部高频约束。

\subsection{训练、权重选定与完整画面推理}

方法3使用 batch size 1、固定学习率 $10^{-4}$ 的 Adam，以FP32完成总计5000步训练。训练从第200步恢复至第5000步；总步数不额外累加。EMA衰减为0.999，逐优化步更新。最终固定第5000步 EMA，后续探索和多数据集测试均使用该权重作为默认模型。

完整画面推理采用8帧时间窗口、时间步长4，空间 tile 为 $64\times64$ packed、步长56。空间和时间重叠处等权平均；不足8帧时重复末帧补齐。输出归一化裁剪后，按相应评测管线恢复 DN 值。分块批大小仅影响一次并行处理多少 tile，不改变训练 batch size。

\subsection{sc635hai 专用可选输出头}
\label{sec:head}

方法8在方法3 EMA基础上冻结主干，仅更新 \texttt{norm}、\texttt{conv\_after\_body} 和 \texttt{conv\_last}，共6个张量、12,400个参数。对应头部由108维 LayerNorm、$108\rightarrow84$线性映射及 $84\rightarrow4$、核为 $1\times3\times3$ 的卷积组成。

适配使用 sc635hai 131x低色温序列，奇偶输入与参考帧池互斥；以0.5概率混入 MIS20S1 回放，减少原域漂移。batch size为1、学习率 $10^{-5}$、训练500步，采用BF16 autocast、FP32损失和0.999 EMA；梯度范数裁剪为1.0。实际500步中，sc635hai采样238步，回放262步。16x运动彩色序列完整保留，不参与反向传播。

该输出头通过既定自动检查并被登记为 sc635hai 专用可选权重；默认方法3继续用于统一比较。下节分别展示默认权重的零样本结果和专用头的适配结果。

\section{多数据集泛化效果展示}
\label{sec:results}

\subsection{展示范围与分组}

本节展示的是\textbf{公司多数据集、跨传感器及跨拍摄条件的泛化表现}。默认方法3只在 MIS20S1暗房真实样本及其合成噪声上训练；SC465SL、sc635hai和sc450ai均以同一权重先做零样本测试。MIS20S1室外用于检查场景变化，1000x用于检查未参与训练的真实增益。另列的方法8结果只表示 sc635hai 专用适配。

45条原始流均完成输入与方法3的固定 ISP 对比。每流取首、中、尾3个连续8帧片段，共24帧，总计1080帧。27条场景流用于观察画质，18条标定流用于噪声与数值诊断；45/45表示推理覆盖。以下按具体数据集给出静态指标和代表画面。

\subsection{MIS20S1：训练覆盖增益与未见真实增益}

\begin{table}[!htbp]
\centering\small
\caption{MIS20S1静态结果。645x在训练覆盖内，使用10x跨增益参考；1000x真实序列未参与训练，使用原生静态参考。}
\label{tab:mis20s1}
\begin{tabular}{lrrrrrr}
\toprule
协议 & 输入PSNR & 方法3 PSNR & 提升 & 输入RMSE & 方法3 RMSE & HF比 \\
 & (dB) & (dB) & (dB) & (DN) & (DN) & \\
\midrule
645x跨增益参考 & 24.3519 & 29.5729 & +5.2210 & 248.1197 & 136.0228 & 0.9712 \\
1000x原生参考 & 44.6204 & 53.2675 & +8.6471 & 24.0565 & 8.8895 & 0.8562 \\
\bottomrule
\end{tabular}
\end{table}

645x的噪声显著下降，高频比0.9712，说明静态图卡结构与参考较为接近。图\ref{fig:mis645}展示同帧输入与最终输出：墙面和图卡周围的彩噪减少，棋盘、文字及卡片轮廓仍可分辨。由于645x参与训练，这一行用于说明训练域的恢复效果，不作为未见增益成绩。

\begin{figure}[!htbp]
\centering\pairhead
\includegraphics[width=\textwidth,trim=0 0 0 44,clip]{fig04_a3_645x_isp.jpg}
\caption{MIS20S1暗房645x固定 ISP 对比，上为全景，下为固定中心裁剪。该图展示最终方法3在训练覆盖域的恢复效果；静态定量只统计有效区域。}
\label{fig:mis645}
\end{figure}

1000x原生参考评测的PSNR提升8.65\db，RMSE由24.06\dn{}降至8.89\dn，图\ref{fig:mis1000}给出对应画面。为检查同流参考的影响，另使用互斥奇偶帧池验证，方法3得到53.10/53.07\db。该结果与完整原生参考评测接近，支持模型对未参与训练的1000x真实流具有静态去噪能力；它仍是同一暗房场景下的真实增益测试。

\begin{figure}[!htbp]
\centering
\makebox[0.32\textwidth]{输入}\hfill\makebox[0.32\textwidth]{方法3（EMA）}\hfill\makebox[0.32\textwidth]{1000x原生参考}\\[0.25em]
\includegraphics[width=\textwidth,trim=0 133 0 135,clip]{fig20_native_1000x.jpg}
\caption{MIS20S1 1000x输入、方法3与原生参考。该序列不参与主方法训练；展示保持相同亮度映射，定量按表\ref{tab:mis20s1}所列协议计算。}
\label{fig:mis1000}
\end{figure}

\subsection{SC465SL与sc635hai：跨传感器零样本静态泛化}

表\ref{tab:zero_shot}中所有结果使用同一方法3 EMA，没有目标传感器训练。SC465SL两档增益的提升为10.52--13.21\db，sc635hai低色温静态提升约13.59\db。两个奇偶方向的结果接近，且高频保持率为0.911--0.958；输出 legal RAW fraction 均为100\%。图\ref{fig:zero_shot_plot}汇总了PSNR提升与高频保持率。

\begin{table}[!htbp]
\centering\small
\caption{不同数据集上的零样本静态泛化。偶/奇表示输入帧池，参考使用相反奇偶池；HF为输出/参考高频 RMS，MAD为时序平均绝对差。}
\label{tab:zero_shot}
\begin{tabular}{llrrrrr}
\toprule
数据集 & 输入池 & 输入PSNR & 方法3 PSNR & 提升 & HF比 & MAD 输入$\to$输出 \\
 & & (dB) & (dB) & (dB) & & (DN) \\
\midrule
SC465SL 256x & 偶 & 41.840 & 55.052 & +13.211 & 0.957 & 31.889$\to$1.567 \\
 & 奇 & 41.840 & 54.999 & +13.158 & 0.958 & 31.886$\to$1.578 \\
SC465SL 2000x & 偶 & 46.039 & 56.563 & +10.523 & 0.911 & 20.238$\to$1.312 \\
 & 奇 & 46.034 & 56.647 & +10.613 & 0.912 & 20.248$\to$1.301 \\
sc635hai 131x & 偶 & 38.737 & 52.325 & +13.588 & 0.919 & 47.481$\to$2.165 \\
 & 奇 & 38.736 & 52.324 & +13.587 & 0.921 & 47.481$\to$2.163 \\
\bottomrule
\end{tabular}
\end{table}

\begin{figure}[!htbp]
\centering
\includegraphics[width=0.85\textwidth]{generalization_static.pdf}
\caption{方法3在三组未参与训练的目标域静态数据上的泛化：左为相对输入的 PSNR 提升，右为高频保持率。数据来自表\ref{tab:zero_shot}，参考均采用互斥奇偶帧池。}
\label{fig:zero_shot_plot}
\end{figure}

SC465SL两档结果在图\ref{fig:sc465_256}和图\ref{fig:sc465_2000}中展示。输入中的细碎彩噪和背景波动得到抑制，风扇、人偶、线缆及背景卡片仍可辨。两图来自完整画面预览，静态定量则按各自的有效区域计算。

\begin{figure}[!htbp]
\centering\pairhead
\includegraphics[width=\textwidth,trim=0 0 0 47,clip]{generalization_sc465sl_256x.png}
\caption{SC465SL暗室256x的零样本结果，上为全景，下为中心裁剪。模型未在SC465SL训练；固定 ISP 下可观察到明显降噪及背景结构保留。}
\label{fig:sc465_256}
\end{figure}

\begin{figure}[!htbp]
\centering\pairhead
\includegraphics[width=\textwidth,trim=0 0 0 47,clip]{generalization_sc465sl_2000x.png}
\caption{SC465SL暗室2000x的零样本结果，使用与256x相同的方法3权重和推理流程。静态PSNR提升约10.52--10.61\db，高频比约0.911--0.912。}
\label{fig:sc465_2000}
\end{figure}

sc635hai 131x在低色温和较大蓝通道白平衡补偿下，仍显示较强的静态去噪能力。图\ref{fig:sc635_static}给出输入、方法3与相反奇偶池静态参考，展示背景卡片和人偶结构的保留。表\ref{tab:zero_shot}对应方法3的零样本结果；后续使用131x训练专用头所得的数字单独列在第\ref{sec:head_results}节，保持两个实验的含义独立。

\begin{figure}[!htbp]
\centering
\makebox[0.32\textwidth]{输入}\hfill\makebox[0.32\textwidth]{方法3（零样本）}\hfill\makebox[0.32\textwidth]{相反奇偶池参考}\\[0.25em]
\includegraphics[width=\textwidth,trim=0 710 0 40,clip]{generalization_sc635hai_131x.png}
\caption{sc635hai 131x低色温数据集的零样本静态对比，选自偶输入方向的第8帧。参考为对侧奇数帧池的静态均值，动态像素不计入静态指标。}
\label{fig:sc635_static}
\end{figure}

\subsection{sc635hai专用输出头：适配域与动态保留集}
\label{sec:head_results}

方法8的131x指标来自训练数据尾段复评，该尾段参与过训练采样，不计为独立泛化结果。表\ref{tab:head_results}显示，专用头在静态域仅带来约0.034\db{}的小幅变化，高频与时序指标保持接近。其16x动态保留集未参加训练；由于缺少干净动态目标，只报告诊断量。

\begin{table}[!htbp]
\centering\small
\caption{sc635hai专用头效果。131x为训练数据尾段复评，16x为完整保留的动态序列；本表与零样本表采用不同评测样本，基线数值不混用。}
\label{tab:head_results}
\begin{tabular}{llrr}
\toprule
数据与协议 & 指标 & 方法3 & 方法8专用头 \\
\midrule
131x适配域 & 静态PSNR，偶方向(dB) & 52.3535 & 52.3875 \\
 & 静态PSNR，奇方向(dB) & 52.3552 & 52.3884 \\
 & HF比（约） & 0.9199 & 0.9181 \\
 & temporal MAD(DN) & 2.1628 & 2.1247 \\
\midrule
16x动态保留集 & HF诊断比 & 0.763605 & 0.763078 \\
 & motion/input ratio & 0.599602 & 0.596438 \\
 & temporal MAD(DN) & 14.6745 & 14.6218 \\
 & 端点裁剪 & 2.1994\% & 1.9344\% \\
 & legal RAW fraction & 100\% & 100\% \\
\bottomrule
\end{tabular}
\end{table}

\begin{figure}[!htbp]
\centering\triplehead
\includegraphics[width=\textwidth,trim=0 0 0 44,clip]{fig11_a8_sc635hai.jpg}
\caption{sc635hai 16x动态保留集：输入、默认方法3和方法8专用输出头。两种输出在结构与纹理上接近，固定 ISP 首、中、末帧复核未见明显新增重影、断边或色偏。}
\label{fig:sc635_head}
\end{figure}

图\ref{fig:sc635_head}中，方法8的变化主要集中在输出映射，未形成可辨的动态结构提升，因此保留其“sc635hai专用可选头”的身份。跨传感器静态泛化的主要证据仍是未经目标域适配的方法3。

\subsection{室外、高增益和sc450ai数据集的实际输出}

多数据集展示也包含强噪声、运动和复杂纹理。MIS20S1室外1500x与训练暗房具有场景和噪声差异。图\ref{fig:outdoor}中方法3明显压低彩噪，但人物局部有姿态平均，树木纹理偏软。该序列没有干净动态参考，因此不提供动态PSNR。

\begin{figure}[!htbp]
\centering\pairhead
\includegraphics[width=\textwidth,trim=0 0 0 44,clip]{fig15_a3_1500x_isp.jpg}
\caption{MIS20S1室外1500x跨场景输出。上为全景，下为中心裁剪；同一默认方法3显著降噪，同时可见人物与背景细节的平滑现象。}
\label{fig:outdoor}
\end{figure}

sc450ai 1435x黑物运动在零样本诊断中，高频比为0.127，端点裁剪约3.28\%，运动/静态时序比为0.057/0.088。图\ref{fig:sc450}保留了方法3原始结果：黑衣内部结构偏空、边界偏软。804x灯光变化的高频比为0.690，暗部仍有低频色斑。SC465SL difficultpattern视频也观察到拖影和重影。上述画面与前三组静态定量共同构成了多数据集展示的完整结果。

\begin{figure}[!htbp]
\centering\pairhead
\includegraphics[width=\textwidth,trim=0 0 600 44,clip]{fig12_a8_sc450ai.jpg}
\caption{sc450ai 1435x黑物运动的零样本结果。此处仅显示原对比图的输入与最终方法3两列；黑色衣物内部与运动轮廓仍偏软。}
\label{fig:sc450}
\end{figure}

\FloatBarrier
表\ref{tab:result_scope}汇总上述数据集的评测协议与实际观察。

{\small
\begin{longtable}{P{0.26\textwidth}P{0.23\textwidth}P{0.40\textwidth}}
\caption{最终方法在不同公司数据集中的展示结果。}\label{tab:result_scope}\\
\toprule
数据条件 & 使用权重与协议 & 已观察结果 \\
\midrule
\endfirsthead
\toprule
数据条件 & 使用权重与协议 & 已观察结果（续） \\
\midrule
\endhead
\bottomrule
\endfoot
MIS20S1暗房645x & 方法3；训练覆盖域、跨增益参考 & +5.22\db，HF 0.971 \\
MIS20S1暗房1000x & 方法3；真实增益留出、原生参考 & +8.65\db；互斥奇偶验证53.10/53.07\db \\
SC465SL暗室256x/2000x & 方法3；跨传感器零样本静态 & +10.52--13.21\db，HF 0.911--0.958 \\
sc635hai低色温131x & 方法3；跨传感器零样本静态 & 约+13.59\db，HF 0.919--0.921 \\
sc635hai动态16x & 方法3及方法8专用头；动态保留集 & 两者画面接近；专用头未见明显新增伪影 \\
MIS20S1室外1500x & 方法3；跨场景动态观察 & 噪声下降，人物姿态与树木结构偏平滑 \\
sc450ai、复杂纹理动态 & 方法3；零样本压力测试 & 可见结构平滑、色斑或重影，不计为通过的静态泛化结果 \\
\end{longtable}
}

\FloatBarrier

\begin{thebibliography}{99}
%========================================================

\bibitem{r1} H. Yue, C. Cao, L. Liao, R. Chu, and J. Yang. Supervised Raw Video Denoising With a Benchmark Dataset on Dynamic Scenes. In \emph{Proceedings of the IEEE/CVF Conference on Computer Vision and Pattern Recognition (CVPR)}, pp. 2301--2310, 2020.

\bibitem{r2} H. Yue, C. Cao, L. Liao, and J. Yang. RViDeformer: Efficient Raw Video Denoising Transformer With a Larger Benchmark Dataset. \emph{IEEE Transactions on Circuits and Systems for Video Technology}, 35(9):8929--8944, 2025. DOI: 10.1109/TCSVT.2025.3553160. 本报告结构图选自预印本v2：\url{https://arxiv.org/abs/2305.00767v2}。

\bibitem{r3} J. Liang, Y. Fan, X. Xiang, R. Ranjan, E. Ilg, S. Green, J. Cao, K. Zhang, R. Timofte, and L. Van Gool. Recurrent Video Restoration Transformer with Guided Deformable Attention. In \emph{Advances in Neural Information Processing Systems (NeurIPS)}, 2022.

\bibitem{r4} J. Lehtinen, J. Munkberg, J. Hasselgren, S. Laine, T. Karras, M. Aittala, and T. Aila. Noise2Noise: Learning Image Restoration without Clean Data. In \emph{Proceedings of the International Conference on Machine Learning (ICML)}, PMLR 80:2965--2974, 2018.

\bibitem{r5} A. Krull, T.-O. Buchholz, and F. Jug. Noise2Void -- Learning Denoising from Single Noisy Images. In \emph{Proceedings of the IEEE/CVF Conference on Computer Vision and Pattern Recognition (CVPR)}, pp. 2129--2137, 2019.

\bibitem{r6} J. Batson and L. Royer. Noise2Self: Blind Denoising by Self-Supervision. In \emph{Proceedings of the International Conference on Machine Learning (ICML)}, PMLR 97:524--533, 2019.

\bibitem{r7} A. Foi, M. Trimeche, V. Katkovnik, and K. Egiazarian. Practical Poissonian-Gaussian Noise Modeling and Fitting for Single-Image Raw-Data. \emph{IEEE Transactions on Image Processing}, 17(10):1737--1754, 2008.

\bibitem{r8} T. Brooks, B. Mildenhall, J. T. Barron, J. Chen, D. Sharlet, and R. Ng. Unprocessing Images for Learned Raw Denoising. In \emph{Proceedings of the IEEE/CVF Conference on Computer Vision and Pattern Recognition (CVPR)}, pp. 11036--11045, 2019.

\bibitem{r9} C. Godard, K. Matzen, and M. Uyttendaele. Deep Burst Denoising. In \emph{Proceedings of the European Conference on Computer Vision (ECCV)}, 2018.

\bibitem{r10} S. W. Hasinoff, D. Sharlet, R. Geiss, A. Adams, J. T. Barron, F. Kainz, J. Chen, and M. Levoy. Burst Photography for High Dynamic Range and Low-Light Imaging on Mobile Cameras. \emph{ACM Transactions on Graphics}, 35(6):192, 2016.

\bibitem{r11} T. Kroeger, R. Timofte, D. Dai, and L. Van Gool. Fast Optical Flow Using Dense Inverse Search. In \emph{Proceedings of the European Conference on Computer Vision (ECCV)}, pp. 471--488, 2016.

\bibitem{r12} E. Reinhard, M. Stark, P. Shirley, and J. Ferwerda. Photographic Tone Reproduction for Digital Images. \emph{ACM Transactions on Graphics}, 21(3):267--276, 2002.

\bibitem{r13} P. Charbonnier, L. Blanc-Feraud, G. Aubert, and M. Barlaud. Two Deterministic Half-Quadratic Regularization Algorithms for Computed Imaging. In \emph{Proceedings of the IEEE International Conference on Image Processing (ICIP)}, vol. 2, pp. 168--172, 1994.

\bibitem{r14} K. Zhang, W. Zuo, Y. Chen, D. Meng, and L. Zhang. Beyond a Gaussian Denoiser: Residual Learning of Deep CNN for Image Denoising. \emph{IEEE Transactions on Image Processing}, 26(7):3142--3155, 2017.

\bibitem{r15} D. Maggioni, E. Sanchez-Monge, and A. Foi. Efficient Multi-Stage Video Denoising With Recurrent Spatio-Temporal Fusion. In \emph{Proceedings of the IEEE/CVF Conference on Computer Vision and Pattern Recognition (CVPR)}, 2021.

\bibitem{r16} M. Tassano, J. Delon, and T. Veit. FastDVDNet: Towards Real-Time Deep Video Denoising Without Flow Estimation. In \emph{Proceedings of the IEEE/CVF Conference on Computer Vision and Pattern Recognition (CVPR)}, 2020.

\bibitem{r17} B. Mildenhall, J. T. Barron, J. Chen, D. Sharlet, R. Ng, and R. Carroll. Burst Denoising with Kernel Prediction Networks. In \emph{Proceedings of the IEEE/CVF Conference on Computer Vision and Pattern Recognition (CVPR)}, 2018.

\bibitem{r18} A. Buades, B. Coll, and J.-M. Morel. A Non-Local Algorithm for Image Denoising. In \emph{Proceedings of the IEEE/CVF Conference on Computer Vision and Pattern Recognition (CVPR)}, 2005.

\bibitem{r19} K. Dabov, A. Foi, V. Katkovnik, and K. Egiazarian. Image Denoising by Sparse 3-D Transform-Domain Collaborative Filtering. \emph{IEEE Transactions on Image Processing}, 16(8):2080--2095, 2007.

\bibitem{r20} S. Guo, Z. Yan, K. Zhang, L. Zhang, and W. Zuo. Toward Convolutional Blind Denoising of Real Photographs. In \emph{Proceedings of the IEEE/CVF Conference on Computer Vision and Pattern Recognition (CVPR)}, 2019.

\end{thebibliography}

\FloatBarrier
\appendix
\section{采纳方法的超参数与模型配置}
\label{app:configuration}

本附录记录最终方法3及方法8的 sc635hai 专用可选头。训练设置以最终 checkpoint 中保存的配置和优化器状态为准，模型与损失以对应实现核对；展示视频的编码以实际产物为准。下列路径均相对于项目根目录。

\subsection{权重与运行环境}

\begin{description}\raggedright
\item[项目目录] \path{/Data4/Xiao/RECOVERED_Stage2_DenoiseModel_20261010}
\item[Python] \path{/home/jingzeliang/miniconda3/envs/aaa_312/bin/python}
\item[方法3最终权重] \path{artifacts/finetune_mis20s1/training/checkpoint_step_005000.pth}，读取 \texttt{ema} 字段。
\item[方法3 SHA-256] {\small\path{97c4332f685f98959b243e4b100434469f714aed6329e100d22b6d8552592e0b}}
\item[方法8专用权重] \path{artifacts/a8_sensor_adaptation/sc635hai/training_20261013/checkpoint_step_000500.pth}，读取 \texttt{ema} 字段。
\item[方法8 SHA-256] {\small\path{8e4980a43767b5b83ecc064e8cdce2f685ea452e94af48a1f4ca252004765b3f}}
\end{description}

\subsection{方法3训练超参数}

\begin{longtable}{P{0.27\textwidth}P{0.65\textwidth}}
\caption{最终主方法的训练与采样配置。}\label{tab:training_config}\\
\toprule 项目 & 实际设置 \\ \midrule
\endfirsthead
\toprule 项目 & 实际设置（续） \\ \midrule
\endhead
\bottomrule\endfoot
初始化 & 官方 \texttt{RViDeformer-L\_CRVD.pth}，strict加载；恢复后的本地预训练文件位于 \path{artifacts/pretrained/} \\
可训练范围 & 全模型，共7,505,716个参数、1208个参数张量 \\
总优化步数 & 5000；从step 200恢复至step 5000 \\
训练batch size & 1；每次优化包含1个8帧样本 \\
学习率 & 固定 $1\times10^{-4}$，无学习率调度器 \\
优化器 & Adam，$\beta_1=0.9$、$\beta_2=0.999$、$\epsilon=10^{-8}$ \\
优化器其他项 & weight decay为0，amsgrad为False \\
数值精度 & FP32；无autocast、GradScaler或梯度累积 \\
梯度裁剪 & 未使用 \\
EMA & decay=0.999，每次optimizer step后更新参数，buffer直接复制；推理读取EMA \\
随机种子 & 20260921；设置Python、NumPy、PyTorch和CUDA随机种子 \\
输入/输出 & $(B,T,C,H,W)=(1,8,4,64,64)$，输出对应8帧 \\
真实:合成比例 & 配置0.5:0.5；512项采样序列的前256项为真实、后256项为合成 \\
真实训练序列 & MIS20S1暗房128x、313x、632x、645x；不含1000x真实序列 \\
真实时间采样 & 随机起点的连续8帧，frame stride=1 \\
合成时间采样 & 每个增益8个噪声变体，从变体池循环取8帧，干净内容相同 \\
裁剪与增强 & 随机 $64\times64$ packed裁剪、真实时窗、合成噪声；未加入翻转或旋转 \\
样本随机状态 & 索引 $i$ 使用 \texttt{seed + i * 104729}；合成变体使用 \texttt{seed + record\_index * 1009 + variant} \\
DataLoader & 默认512项，shuffle=False，num\_workers=0，drop\_last=True；耗尽后重新遍历 \\
训练输出裁剪 & 预测先clamp到 $[0,1]$，再计算损失 \\
最终权重选择 & step-5000 EMA；未将后续运动、统一微调或门控候选合并进去 \\
\end{longtable}

实际配置文件为 \path{artifacts/finetune_mis20s1/training/a3_config.json}；训练及采样实现为 \path{src/raw_fusion/finetune/a3.py}。

\subsection{参考、归一化与损失设置}

\begin{longtable}{P{0.27\textwidth}P{0.65\textwidth}}
\caption{主方法监督与损失参数。}\label{tab:loss_config}\\
\toprule 项目 & 设置 \\ \midrule
\endfirsthead
\toprule 项目 & 设置（续） \\ \midrule
\endhead
\bottomrule\endfoot
主参考 & MIS20S1暗房10x的200帧；packed形状 $(4,540,960)$ \\
通道顺序 & R、Gr、Gb、B \\
主参考黑电平 & R=266.125、Gr=275.625、Gb=275.375、B=267.0\dn \\
白电平 & 4095\dn \\
真实输入映射 & 按每条流manifest中的整数位移配准；逐通道slope/intercept映射到参考域，再按式\eqref{eq:normalization}归一化 \\
有效掩码 & 参考有效区 $\cap$ 8帧均静态 $\cap$ 全部8帧/4通道均未饱和；合成样本不另加运动掩码 \\
饱和阈值 & 映射参考域3800\dn，归一化阈值按各参考通道黑电平换算 \\
运动阈值 & 对应采样位置的运动mask $<0.5$，全部8帧同时满足 \\
损失总式 & $\mathcal L_{\mathrm{char}}+0.1\mathcal L_{\mathrm{grad}}+0.1\mathcal L_{\mathrm{hp}}$ \\
Charbonnier & $\sqrt{(\hat y-y)^2+10^{-6}}$；标准 $\epsilon^2$ 记法下 $\epsilon=10^{-3}$ \\
梯度损失 & 水平、垂直梯度幅值差的masked L1，两方向各0.5；相邻mask相乘 \\
高通损失 & $H(u)=u-\mathrm{AvgPool3d}(u)$ 后的masked L1 \\
高通池化 & kernel=$(1,3,3)$，stride=1，padding=$(0,1,1)$；默认零填充且包含填充计数 \\
掩码归约 & 掩码沿通道展开；分母为mask.sum().clamp\_min(1)，逐有效元素平均 \\
动态损失/教师蒸馏 & 最终方法3未使用 \\
\end{longtable}

参考统计与空间资源分别见 \path{artifacts/finetune_mis20s1/supervision/references/a1_reference_report.json} 和同目录 \path{a1_reference_manifest.json}。配准、曝光系数与运动掩码逐流记录于 \path{artifacts/finetune_mis20s1/supervision/stream_manifests/}，因此不以单组曝光系数替代所有流。

\subsection{网络构造参数}

\begin{longtable}{P{0.31\textwidth}P{0.61\textwidth}}
\caption{RViDeformer-L在最终方法中的实际构造配置。}\label{tab:model_config}\\
\toprule 参数 & 设置 \\ \midrule
\endfirsthead
\toprule 参数 & 设置（续） \\ \midrule
\endhead
\bottomrule\endfoot
输入/输出通道 & 4 / 4 \\
\texttt{upscale} & 1 \\
\texttt{img\_size} & [8, 64, 64] \\
\texttt{window\_size} & [6, 8, 8] \\
\texttt{depths} & [4, 4, 4, 4, 4, 4, 4, 2, 2, 2, 2] \\
\texttt{embed\_dims} & [84, 84, 84, 84, 84, 84, 84, 108, 108, 108, 108] \\
\texttt{num\_heads} & [6, 6, 6, 6, 6, 6, 6, 6, 6, 6, 6] \\
\texttt{indep\_reconsts} & [11, 12]，按构造实参记录 \\
\texttt{mul\_attn\_ratio} & 0.75 \\
\texttt{mlp\_ratio} & 2 \\
\texttt{qkv\_bias / qk\_scale} & True / None \\
\texttt{drop\_path\_rate} & 0.2（配置最大值） \\
归一化 & LayerNorm \\
\texttt{pa\_frames} & 2 \\
\texttt{deformable\_groups} & 12 \\
\texttt{spynet\_path} & None \\
\texttt{recal\_all\_flows} & False \\
\texttt{nonblind\_denoising} & False \\
激活检查点构造参数 & \texttt{use\_checkpoint\_attn=False}，\texttt{use\_checkpoint\_ffn=False}；\texttt{no\_checkpoint\_attn\_blocks} 和 \texttt{no\_checkpoint\_ffn\_blocks} 均为空列表 \\
前7级空间尺度 & [1, 2, 4, 8, 4, 2, 1]；先降采样再恢复 \\
参数与buffer & 参数7,505,716个；buffer 57个，共2,668,800个元素 \\
状态字典 & 每套model/EMA共1265张量，其中1208个为参数张量；张量个数不等于参数量 \\
\end{longtable}

模型配置入口为 \path{src/raw_fusion/finetune/a3.py} 的 \texttt{build\_rvideformer}，结构实现为 \path{src/raw_fusion/finetune/rvideformer_arch.py}。

\subsection{噪声合成参数}

表\ref{tab:shot}和表\ref{tab:read}列出实际生成器使用的逐档经验参数。增益标签 $g$ 对应的 $G$ 为标签除以100；shot系数已包含该档增益，不应在式\eqref{eq:synthetic}中再乘一次 $G$。表中数值保留6位小数，精确复现读取原JSON的完整精度。

\begin{table}[!htbp]
\centering\small
\caption{逐增益经验shot系数 $\alpha_{c,g}$，单位DN。}
\label{tab:shot}
\begin{tabular}{rrrrrr}
\toprule
增益标签 & $G$ & R & Gr & Gb & B \\
\midrule
100 & 1 & 0.146240 & 0.133124 & 0.134059 & 0.153964 \\
200 & 2 & 0.303604 & 0.262204 & 0.263975 & 0.301672 \\
400 & 4 & 0.595802 & 0.525547 & 0.529092 & 0.569351 \\
800 & 8 & 1.234060 & 1.038338 & 1.045279 & 1.185013 \\
1600 & 16 & 2.553175 & 2.271041 & 2.286595 & 2.535520 \\
3200 & 32 & 5.057363 & 4.411844 & 4.442612 & 4.909565 \\
6400 & 64 & 10.160087 & 8.932881 & 8.913312 & 10.015178 \\
12800 & 128 & 24.593596 & 24.830152 & 24.722971 & 25.129168 \\
25600 & 256 & 48.717270 & 50.454506 & 50.209071 & 50.189296 \\
\bottomrule
\end{tabular}
\end{table}

\begin{table}[!htbp]
\centering\small
\caption{逐增益read variance $v_{c,g}$，单位DN$^2$。}
\label{tab:read}
\begin{tabular}{rrrrrr}
\toprule
增益标签 & $G$ & R & Gr & Gb & B \\
\midrule
100 & 1 & 1.582505 & 1.582505 & 1.582505 & 1.582505 \\
200 & 2 & 2.813700 & 2.813700 & 2.813700 & 2.813700 \\
400 & 4 & 8.617427 & 8.616957 & 8.616957 & 8.617427 \\
800 & 8 & 29.718838 & 29.718838 & 29.718838 & 29.718838 \\
1600 & 16 & 4.396071 & 4.396071 & 4.396071 & 4.396071 \\
3200 & 32 & 11.254265 & 11.254265 & 11.254265 & 11.254265 \\
6400 & 64 & 39.564642 & 34.466890 & 34.466890 & 34.467829 \\
12800 & 128 & 114.346263 & 114.346263 & 114.346263 & 114.346263 \\
25600 & 256 & 364.052709 & 364.046605 & 364.049657 & 364.049657 \\
\bottomrule
\end{tabular}
\end{table}

生成器使用的聚合通道相关矩阵为
\begin{equation}
C=\begin{bmatrix}
1 & 0.01588756 & 0.00647320 & 0.00093781\\
0.01588756 & 1 & 0.00753526 & 0.00049242\\
0.00647320 & 0.00753526 & 1 & 0.02030878\\
0.00093781 & 0.00049242 & 0.02030878 & 1
\end{bmatrix}.
\end{equation}
固定图样偏置读取主参考manifest中的 \texttt{fpn\_paths.offset\_dn} 空间数组。每档8个噪声变体使用独立随机数，随后先round至1\dn{}格点、再clip至 $[0,4095]$，保存为uint16。完整模型位于 \path{artifacts/finetune_mis20s1/noise_model/noise_model.json}；生成器位于 \path{src/raw_fusion/finetune/noise.py} 的 \texttt{generate\_noisy\_frame}。

\subsection{推理、评测与固定 ISP 参数}

\begin{longtable}{P{0.27\textwidth}P{0.65\textwidth}}
\caption{最终主方法的推理与展示配置。}\label{tab:inference_config}\\
\toprule 项目 & 设置 \\ \midrule
\endfirsthead
\toprule 项目 & 设置（续） \\ \midrule
\endhead
\bottomrule\endfoot
模型状态 & EMA，eval模式；不更新参数 \\
连续视频时间窗口 & 8帧，窗口步长4；不足8帧时重复末帧，丢弃补齐位置的输出 \\
空间tile/步长 & $64\times64$ packed / 56 packed；完整画面覆盖 \\
重叠区域 & 时域和空域均采用等权累计、除以覆盖计数 \\
分块并行批大小 & 全数据集展示使用tile\_batch=12；跨传感器审计入口配置为2，原始全帧评估入口常用8；均不改变训练bs=1 \\
输出裁剪 & 归一化域 $[0,1]$，再按对应管线恢复DN \\
跨传感器静态采样 & 输入与参考奇偶帧池互斥；每片段8帧、原视频stride=2。SC465SL每方向5片段（40帧），sc635hai零样本每方向3片段（24帧），专用头尾段复评每方向1片段（8帧） \\
静态有效区 & 奇偶参考置信度 $\ge0.9$ \\
展示采样 & 每流首、中、尾各1个连续8帧片段，共24帧；45流共1080帧 \\
去马赛克 & canonical Bayer重建，OpenCV edge-aware \\
曝光选择 & 每条流由输入99.5百分位确定；Reinhard后线性目标0.8，EV范围 $[-6,8]$；输入与输出共享冻结值 \\
白平衡 & 各流manifest中的固定WB，输入与输出共享 \\
色调与高光 & Reinhard tone mapping；高光重建开启；随后sRGB gamma \\
其他ISP项 & 2DNR/3DNR关闭，false color suppression关闭 \\
色彩模式 & uncalibrated camera RGB；未使用传感器ColorMatrix \\
实际预览视频 & HEVC（libx265），CRF16，yuv444p，1920$\times$1080，6fps；每视频24帧 \\
面板布局 & 输入左、方法3右；上为全景、下为固定中心裁剪 \\
\end{longtable}

表\ref{tab:inference_config}的 ISP 与视频项对应45流全数据集预览。奇偶参考面板和方法8对比使用原审计固定渲染：固定白平衡与曝光、普通 Bayer 去马赛克、sRGB变换，不包含 Reinhard 与高光重建；各对比内部共享参数。

分块实现见 \path{src/raw_fusion/finetune/a4.py}；跨传感器清单见 \path{configs/a3_cross_sensor_zero_shot_20261012.json} 和 \path{configs/a3_cross_sensor_extension_20261012.json}；全数据集设置见 \path{configs/a3_all_dataset_traditional_isp_preview_20261017.json}。该配置的 \texttt{preferred\_encoder} 记为libx264，但45个实际视频及对应 \texttt{video\_manifest.json} 均为libx265，表\ref{tab:inference_config}按实际产物记录。

\subsection{方法8的sc635hai专用头完整配置}

\begin{longtable}{P{0.27\textwidth}P{0.65\textwidth}}
\caption{专用可选输出头的训练设置。}\label{tab:head_config}\\
\toprule 项目 & 实际设置 \\ \midrule
\endfirsthead
\toprule 项目 & 实际设置（续） \\ \midrule
\endhead
\bottomrule\endfoot
初始化与范围 & 方法3 step-5000 EMA；主干冻结且eval，仅训练输出头 \\
可训练模块 & \texttt{norm}、\texttt{conv\_after\_body}、\texttt{conv\_last}，共6张量、12,400参数 \\
头部结构与参数 & LayerNorm108：216；Linear108$\to$84：9156；Conv3d84$\to$4，kernel=$(1,3,3)$：3028 \\
目标传感器样本 & sc635hai 131x低色温，共150帧，完整范围 $[0,150)$ \\
输入/目标帧池 & 偶池0,2,\ldots,148与奇池1,3,\ldots,149互换；每池75帧，目标为对侧静态均值 \\
时间窗口 & 同一池8帧，即原视频frame stride=2 \\
训练batch / 步数 & 1 / 500 \\
随机种子 & 20261013 \\
学习率与优化器 & 固定 $10^{-5}$，Adam，betas=$(0.9,0.999)$，eps=$10^{-8}$，weight decay=0，amsgrad=False；无调度器 \\
精度与损失 & BF16 autocast，损失转FP32；与式\eqref{eq:loss}相同，权重1/0.1/0.1 \\
EMA / 梯度裁剪 & decay=0.999 / gradient norm=1.0 \\
patch与有效区域 & $64\times64$ packed；静态置信度 $\ge0.9$，输入未饱和；有效比例至少0.25，最多尝试32次裁剪 \\
目标域增强 & 水平翻转、垂直翻转、转置，各以0.5概率施加 \\
回放概率与实绩 & MIS20S1回放概率0.5；实际回放262步，sc635hai采样238步；回放沿用主方法数据管线 \\
检查点间隔 & 每100步保存，最终保留step-500 EMA \\
131x静态复评片段 & 偶输入134,136,\ldots,148；奇输入135,137,\ldots,149。该尾段参与训练采样，不作为独立留出 \\
独立动态保留集 & sc635hai 16x运动彩色200帧，不参与训练；仅无参考诊断与固定ISP观察 \\
用途 & sc635hai专用可选头，不替代默认方法3 \\
\end{longtable}

配置为 \path{configs/a8_sc635hai_adaptation_20261013.json}；实际训练摘要与保存配置位于 \path{artifacts/a8_sensor_adaptation/sc635hai/training_20261013/}；实现为 \path{src/raw_fusion/finetune/a8_adaptation.py}。

\subsection{结果来源与复现入口}

主方法静态评估位于 \path{artifacts/finetune_mis20s1/evaluation/}；跨传感器静态结果位于 \path{artifacts/a3_cross_sensor_zero_shot_20261012/} 和 \path{artifacts/a3_cross_sensor_extension_20261012/}；45流完整预览位于 \path{artifacts/a3_all_dataset_traditional_isp_preview_20261017/}。专用输出头的解释与检查结果见 \path{docs/A8_SC635HAI_ADAPTATION_RESULT_2026-10-13.md}。这些文件名保留原始实验标识，以便定位已有产物。

从项目根目录执行以下只评测命令，可按对应协议复现主要结果；命令中的 \texttt{python} 指本附录记录的Python环境。

{\small
\begin{verbatim}
python scripts/run_a4_eval.py --max-frames 200 \
  --tile-batch 8 \
  --output-dir artifacts/finetune_mis20s1/evaluation
python scripts/run_native_1000x_eval.py --max-frames 200 \
  --tile-batch 8 \
  --output-dir artifacts/finetune_mis20s1/evaluation/native_1000x
python scripts/evaluate_a3_cross_sensor.py \
  --config configs/a3_cross_sensor_zero_shot_20261012.json \
  --output-dir artifacts/a3_cross_sensor_zero_shot_20261012 \
  --phase all --device cuda:0 --tile-batch 2
python scripts/evaluate_a3_cross_sensor.py \
  --config configs/a3_cross_sensor_extension_20261012.json \
  --output-dir artifacts/a3_cross_sensor_extension_20261012 \
  --phase all --device cuda:0 --tile-batch 2
\end{verbatim}
}

结果表中的645x为训练覆盖域复评，1000x为真实增益留出；零样本跨传感器表只使用方法3，sc635hai专用头的131x尾段属于训练数据复评。复现时应沿用这些数据身份与各自参考定义。

\section{补充探索的记录与图表来源}
\label{app:djc_sources}

第\ref{sec:djc}节依据“任务三-djc”文件夹下的 \path{deliverable_20261025/} 整理；以下路径均相对此交付目录。核查范围包括研究报告与台账、服务器历史实验、本地训练与补充验证、模型实现和图像来源。正文按研究问题归纳，不将部署自检、草案预算和实际优化器更新混为一谈。

\begin{longtable}{P{0.23\textwidth}P{0.66\textwidth}}
\caption{ExposureDiffusion路线的主要证据入口。}\label{tab:djc_sources}\\
\toprule
内容 & 对应记录 \\
\midrule
\endfirsthead
\toprule
内容 & 对应记录 \\
\midrule
\endhead
\bottomrule
\endfoot
早期噪声研究 & \path{project/code/noise_model/} 下的 \path{TWO_SHOTS_CVPR2026.md}、\path{ONE_FRAME_DARK_VALIDATION.md}、\path{PNNP_DECOMPOSITION_COMPARISON.md}、\path{QUANTIZED_LIKELIHOOD.md} 与 \path{STATIC_CONDITIONAL_RESIDUAL.md} \\
多传感器自标定 & \path{archive/server/calibration/report.md} \\
单帧1000步适配 & \path{archive/server/reports/gpu_start/completion_audit.json}，及 \path{archive/server/runs/company_mixed_gpu_001/} \\
历史模块试验与解冻 & \path{archive/server/experiments/}、\path{archive/server/runs/}、\path{archive/server/evaluation/}；汇总表 \path{results/tables/historical_validation_all_metrics.csv} \\
最终五帧结构与实际训练 & \path{project/code/company_exposure/burst_detail/model.py}；\path{project/code/scripts/exposure_output_residual_local.py}；\path{project/work_v3/runs/residual_v3_001/} \\
固定与补充验证 & 上述运行目录中的 \path{fixed_validation.json}、\path{baseline_validation.json} 与 \path{validation_2300.json}；\path{results/tables/supplement_validation.csv}、\path{supplement_sensitivity.csv} \\
165例图像与权重身份 & \path{results/report_comparisons/} 中各流的 \path{best_full.png}、\path{best_detail.png} 及 \path{best_provenance.json}；\path{results/all165/inference_contract.json} \\
\end{longtable}

早期噪声建模的数值按上述研究记录转述；对应 \path{runs/noise/} 的原始拟合JSON与NPZ未包含在本次交付包中，故本次仅核对其报告与实现，不将它们作为新完成的独立验收。后续历史训练的352项汇总指标已与原运行JSON逐项比对；五帧候选的固定验证、补充裁块和图像权重身份亦分别核对。

本报告插入的五张场景图保持原始PNG内容，排版仅裁去原小字号表头并重排标签。复制来源、文件哈希和最佳权重身份保存在本项目 \path{docs/assets/stage2_report_20261023/djc/figure_sources.json}。曲线与分域图由同目录 \path{plot_djc_results.py} 基于随附最小CSV重绘，来源记录为 \path{plot_sources.json}。这组图中的官方 RViDeformer 与第\ref{sec:adopted}节的微调方法3始终作为不同权重处理。

\end{document}
